% Лабораторная работа 1. Моделирование и устойчивость. Вариант 3.
% Основа: DCS.pdf, с. 17-22.

\setlength{\abovedisplayskip}{7pt}
\setlength{\belowdisplayskip}{7pt}
\setlength{\abovedisplayshortskip}{4pt}
\setlength{\belowdisplayshortskip}{4pt}

% Рисунки и таблицы располагаются рядом с текстом.
\makeatletter
\newenvironment{reportblock}[1]{%
  \par\addvspace{8pt}\noindent
  \begin{minipage}{\textwidth}
    \def\@captype{#1}%
    \centering
}{%
  \end{minipage}\par\vspace{4pt}
}
\makeatother

\newcommand{\reportimage}[4][0.9\textwidth]{%
  \begin{reportblock}{figure}
    \includegraphics[width=#1,height=0.65\textheight,keepaspectratio]{#2}%
    \caption{#3}
    \label{#4}
  \end{reportblock}
}

\section{Цель работы}

Исследовать влияние дискретизации на устойчивость и характер переходных
процессов, подобрать обратную связь по заданным собственным числам
и построить дискретные командные генераторы.

\section{Задание 1. Влияние дискретного элемента на непрерывную систему}

Для варианта 3 $T=0{,}2$ с, $K_{CO}=4$. На вход подаётся единичное воздействие, $y(0)=0$. Моделирование выполнено на интервале $0\leq t\leq10$ с
с шагом интегрирования $0{,}01$ с.

\reportimage{task1_scheme.png}
  {Схема системы с экстраполятором нулевого порядка}
  {fig:task1-scheme}

\subsection{Устойчивость и выбор коэффициента}

Обозначим $y_k=y(kT)$, $e_k=1-y_k$. При удержании управления
$u(t)=K_{FB}e_k$ на интервале $kT\leq t<(k+1)T$ получаем
\[
  \dot y=K_{CO}u,\qquad
  y_{k+1}=y_k+TK_{CO}K_{FB}(1-y_k).
\]
Следовательно,
\[
  e_{k+1}=qe_k,\qquad q=1-0{,}8K_{FB},\qquad y_k=1-q^k.
\]
Условие асимптотической устойчивости:
\[
  |q|<1\quad\Longleftrightarrow\quad 0<K_{FB}<2{,}5.
\]
При $K_{FB}=0$ имеем нейтральную границу ($q=1$), при
$K_{FB}=2{,}5$ \textemdash{} колебательную ($q=-1$).
Монотонный процесс соответствует $0<K_{FB}\leq1{,}25$,
затухающие колебания \textemdash{} $1{,}25<K_{FB}<2{,}5$.

Наиболее быстрое точное установление достигается при
\[
  q=0,\qquad K_{FB}=\frac{1}{TK_{CO}}=1{,}25:
  \qquad y(t)=1\quad\text{при }t\geq T=0{,}2\text{ с}.
\]
При приближении $K_{FB}$ к $2{,}5$ снизу затухание замедляется.
Максимум колебательности среди асимптотически устойчивых режимов
не достигается; на границе колебания становятся незатухающими.

\begin{reportblock}{table}
  \caption{Исследованные режимы}
  \label{tab:task1-gain}
  \begin{tabularx}{0.95\textwidth}{c c X}
    \toprule
    $K_{FB}$ & $q$ & Результат моделирования \\
    \midrule
    $0$ & $1$ & Выход остаётся нулевым \\
    $2{,}5$ & $-1$ & Незатухающие колебания от 0 до 2 \\
    $2{,}4$ & $-0{,}92$ & Затухающие колебания около 1 \\
    $1$ & $0{,}2$ & Монотонное приближение к 1 \\
    $1{,}25$ & $0$ & Достижение 1 за один период $T$ \\
    \bottomrule
  \end{tabularx}
\end{reportblock}

\reportimage{task1_responses.png}
  {Переходные процессы при пяти значениях $K_{FB}$}
  {fig:task1-responses}

Результаты подтверждают расчёт: при $K_{FB}=2{,}5$ отсчёты
чередуются как $0,2,0,2,\ldots$, период колебаний равен $2T=0{,}4$ с.
Максимальное расхождение отсчётов моделирования и формулы $y_k=1-q^k$
не превышает $3{,}6\cdot10^{-15}$.

Без экстраполятора $\dot e=-4K_{FB}e$, поэтому система устойчива
при любом $K_{FB}>0$. Дискретизация ограничивает усиление сверху:
$K_{FB}<2{,}5$. Быстродействие максимально при $K_{FB}=1{,}25$;
дальнейшее увеличение усиления вызывает колебания и замедляет затухание.

\clearpage
\section{Задание 2. Исследование устойчивости дискретных систем}

\subsection{Модель объекта и дискретизация}

Для объекта $\ddot y=u$ выберем $x=(y,\dot y)^{\mathsf T}$.
При $T=0{,}2$ с и $x(0)=(1,0)^{\mathsf T}$ непрерывная модель имеет вид
\[
  \dot x=A_{\text{н}}x+B_{\text{н}}u,\qquad y=Cx,
  \qquad
  A_{\text{н}}=\begin{pmatrix}0&1\\0&0\end{pmatrix},\quad
  B_{\text{н}}=\begin{pmatrix}0\\1\end{pmatrix},\quad
  C=\begin{pmatrix}1&0\end{pmatrix}.
\]
Поскольку $A_{\text{н}}^2=0$, имеем $e^{A_{\text{н}}\tau}=I+A_{\text{н}}\tau$.
Точная дискретизация с удержанием управления даёт
\[
  A=e^{A_{\text{н}}T}=\begin{pmatrix}1&0{,}2\\0&1\end{pmatrix},\qquad
  B=\int_0^T e^{A_{\text{н}}\tau}B_{\text{н}}\,d\tau
   =\begin{pmatrix}T^2/2\\T\end{pmatrix}
   =\begin{pmatrix}0{,}02\\0{,}2\end{pmatrix}.
\]
Таким образом, $x(k+1)=Ax(k)+Bu(k)$, $y(k)=Cx(k)$.

\subsection{Расчёт обратной связи}

При $u(k)=-Kx(k)$, $K=(k_1\;\;k_2)$ матрица замкнутой системы
и её характеристический полином равны
\[
  F=A-BK=
  \begin{pmatrix}
    1-0{,}02k_1 & 0{,}2-0{,}02k_2\\
    -0{,}2k_1 & 1-0{,}2k_2
  \end{pmatrix},
\]
\[
  \det(zI-F)=z^2+(0{,}02k_1+0{,}2k_2-2)z
                 +1+0{,}02k_1-0{,}2k_2.
\]
Для желаемых корней положим $s=z_1+z_2$, $p=z_1z_2$.
Приравнивая коэффициенты полиному $p^*(z)=z^2-sz+p$, получаем
\[
  k_1=\frac{1-s+p}{T^2},\qquad
  k_2=\frac{3-s-p}{2T}.
\]

\begin{reportblock}{table}
  \caption{Желаемые корни и коэффициенты обратной связи}
  \label{tab:task2-roots}
  \begin{tabular}{c c c r r}
    \toprule
    Набор & Корни $z_1,z_2$ & $p^*(z)$ & $k_1$ & $k_2$ \\
    \midrule
    1 & $0{,}2;\;0{,}8$ & $z^2-z+0{,}16$ & $4$ & $4{,}6$ \\
    2 & $0{,}5;\;-0{,}8$ & $z^2+0{,}3z-0{,}4$ & $22{,}5$ & $9{,}25$ \\
    3 & $0{,}1;\;-0{,}3$ & $z^2+0{,}2z-0{,}03$ & $29{,}25$ & $8{,}075$ \\
    4 & $\pm0{,}2j$ & $z^2+0{,}04$ & $26$ & $7{,}4$ \\
    5 & $-0{,}3\pm0{,}7j$ & $z^2+0{,}6z+0{,}58$ & $54{,}5$ & $7{,}55$ \\
    \bottomrule
  \end{tabular}
\end{reportblock}

Собственные числа матриц $F$ совпали с заданными с точностью
до округления. Для всех пяти наборов $\rho(F)\leq0{,}8<1$,
следовательно, системы асимптотически устойчивы.

\subsection{Результаты моделирования}

Блок состояния выдаёт вектор $x$, по которому вычисляются $u=-Kx$
и $y=Cx$. Для всех наборов $x(0)=(1,0)^{\mathsf T}$,
интервал моделирования $0\leq t\leq10$ с.

\reportimage{task2_scheme.png}
  {Схема дискретной системы с обратной связью по состоянию}
  {fig:task2-scheme}

\reportimage[0.85\textwidth]{task2_responses.png}
  {Свободные движения системы для пяти наборов собственных чисел}
  {fig:task2-responses}

Все процессы стремятся к нулю. Набор 1 даёт монотонный процесс,
набор 2 \textemdash{} слабые колебания около нуля. Наборы 3 и 4
быстро затухают, в наборе 5 колебания выражены сильнее.
Меньший модуль корней ускоряет асимптотическое затухание;
отрицательные и комплексные корни задают колебательные составляющие.

Расхождение с точным решением $y(k)=C(A-BK)^k x(0)$
не превышает $4{,}5\cdot10^{-16}$. Графики согласуются
с корневым критерием устойчивости и рассчитанной обратной связью.

\clearpage
\section{Задание 3. Построение дискретных командных генераторов}

\subsection{Генератор гармонического сигнала}

Для варианта 3 требуется сформировать
\[
  g(k)=\sin(kT\omega),\qquad
  \omega=0{,}08\ \text{рад/с},\qquad T=0{,}2\ \text{с}.
\]
Положим $\theta=\omega T=0{,}016$ и
$\xi_h(k)=(\sin(k\theta),\cos(k\theta))^{\mathsf T}$.
По формулам сложения получаем автономную модель
\[
  \xi_h(k+1)=\Gamma_h\xi_h(k),\qquad g(k)=H_h\xi_h(k),
\]
\[
  \Gamma_h=R(\theta),\qquad
  R(\varphi)=
  \begin{pmatrix}
    \cos\varphi&\sin\varphi\\
    -\sin\varphi&\cos\varphi
  \end{pmatrix},\qquad
  H_h=\begin{pmatrix}1&0\end{pmatrix},\qquad
  \xi_h(0)=\begin{pmatrix}0\\1\end{pmatrix}.
\]
В реализации «вход-состояние-выход» $B_h=0_{2\times1}$, $D_h=0$,
внешний вход равен нулю. Сигнал формируется за счёт начального состояния.
Собственные числа $e^{\pm j\theta}$ имеют единичный модуль,
поэтому амплитуда колебаний сохраняется.

\reportimage[0.8\textwidth]{task3_harmonic_scheme.png}
  {Схема дискретного генератора гармонического сигнала}
  {fig:task3-harmonic-scheme}

\reportimage{task3_harmonic_signal.png}
  {Сравнение выхода генератора и заданного гармонического сигнала}
  {fig:task3-harmonic-signal}

На интервале $0\leq t\leq160$ с получен 801 отсчёт сигнала
с амплитудой 1. Максимальное расхождение с $\sin(0{,}08kT)$
не превышает $4{,}4\cdot10^{-15}$.

\clearpage
\subsection{Генератор внешнего возмущения}

Задан сигнал
\[
  g(k)=2e^{-kT}\sin(3kT)+0{,}1(kT)^2,\qquad T=0{,}25\ \text{с}.
\]
Обозначим $\eta_1(t)=2e^{-t}\sin(3t)$,
$\eta_2(t)=2e^{-t}\cos(3t)$, $p(t)=0{,}1t^2$.
В моменты $t=kT$ выберем состояние
$\xi_d=(\eta_1,\eta_2,p,\dot p,\ddot p)^{\mathsf T}$.
Для гармонической части переход задаётся матрицей $e^{-T}R(3T)$,
а для полинома и его производных \textemdash{} матрицей точного
сдвига $P$, полученной из ряда Тейлора. В автономной модели
$\xi_d(k+1)=\Gamma_d\xi_d(k)$, $g(k)=H_d\xi_d(k)$ имеем
\[
  \Gamma_d=\operatorname{diag}\bigl(e^{-T}R(3T),P\bigr),
  \qquad
  P=\begin{pmatrix}
    1&T&T^2/2\\
    0&1&T\\
    0&0&1
  \end{pmatrix}
  =\begin{pmatrix}
    1&0{,}25&0{,}03125\\
    0&1&0{,}25\\
    0&0&1
  \end{pmatrix},
\]
\[
  H_d=\begin{pmatrix}1&0&1&0&0\end{pmatrix},\qquad
  \xi_d(0)=\begin{pmatrix}0&2&0&0&0{,}2\end{pmatrix}^{\mathsf T},
  \qquad B_d=0_{5\times1},\quad D_d=0.
\]
\reportimage[0.65\textwidth]{task3_disturbance_scheme.png}
  {Схема дискретного генератора внешнего возмущения}
  {fig:task3-disturbance-scheme}

\reportimage[0.85\textwidth]{task3_disturbance_signal.png}
  {Сравнение выхода генератора и заданного возмущения}
  {fig:task3-disturbance-signal}

На интервале $0\leq t\leq10$ с получен 41 отсчёт.
Начальные колебания затухают, затем преобладает рост $0{,}1t^2$.
Максимальное расхождение с заданной функцией не превышает
$1{,}8\cdot10^{-15}$.

\section{Выводы}

Дискретизация ограничивает устойчивость условием $0<K_{FB}<2{,}5$;
наиболее быстрое точное установление получено при $K_{FB}=1{,}25$.
Все пять наборов коэффициентов обратной связи обеспечивают заданные
устойчивые корни и затухание свободных движений. Генераторы воспроизводят
гармонический сигнал и возмущение с расхождениями на уровне
погрешности вычислений.
